\renewcommand{\thesubsection}{\Alph{subsection}}

% \appendix

\section*{Appendix}
\subsection{Additional ablations}
Here we provide additional ablation experiments that supplement the model choices from the main paper.

\paragraph{Support points.}


In addition to global and local support grids described in the paper, we also test using the SIFT detector~\cite{lowe04distinctive} to select support points in \cref{fig:accuracy_global_local_sift}.
While adding SIFT support points is better than sampling points uniformly on a global regular grid, adding a local grid yields significantly better results.
The combination of global and local grids is slightly better, and in  \cref{fig:accuracy_global_local} we further test this combination with different grid sizes ranging from 4 $\times$ 4 to 10 $\times$ 10.
We find that a configuration of 25 global and 64 local points is optimal for TAP-Vid-DAVIS in terms of speed and accuracy.
While there are differences in performance across these hyper-parameters, the overall impact on performance is less than one percentage point for a wide variety of configurations.
We use the selected configuration for evaluation on all the TAP-Vid benchmarks.


\begin{figure}[tb]
    \begin{minipage}{.47\linewidth}
        \centering
        \setlength\tabcolsep{0.9pt}
       
\includegraphics[width=\textwidth]{figures/plot_global_local_sift.pdf}
\caption{\textbf{Accuracy on TAP-Vid-DAVIS} depending on the number of additional local, global grid and SIFT support points. When combining point sources, we use (close to) the same number of points.}%
\label{fig:accuracy_global_local_sift}

    \end{minipage}%
    \hspace{0.02\linewidth}
    \begin{minipage}{0.5\linewidth}
        \centering
        \setlength\tabcolsep{1pt}
        \includegraphics[width=0.9\textwidth]{figures/plot_global_local.pdf}
\caption{\textbf{Number of support points.} Accuracy for additional global grid and local grid points. We compute $\delta_{avg}^{vis}$ on TAP-Vid-DAVIS to find the optimal grid sizes for TAP-Vid benchmarks. We use 25 \& 64 because of speed and higher occlusion accuracy ($89.3$ vs. $88.7$).}%
\label{fig:accuracy_global_local}

    \end{minipage}
\end{figure}

\paragraph{Training sliding window size.}

In, \cref{tab:tab_train_sliding_window_len_ablation} we ablate the sliding window size during training.
While \method can benefit from larger window sizes, the length of our training sequences from Kubric~\cite{tapvid} is limited to 24 frames.
Therefore, there is a trade-off between the context length and the ability to propagate predictions to future sliding windows.
A sliding window of 8 is optimal for our training dataset.


\paragraph{Inference sliding window size.}

In \cref{tab:tab_inf_sliding_window_len_ablation}, we examine the impact of the sliding window size on the performance of a model trained with a sliding window of $T=8$.
We find that the model performs best when the training and evaluation window sizes are identical.

\begin{table}[tb]
    \begin{minipage}{.47\linewidth}
        \centering
        \setlength\tabcolsep{0.9pt}
       %\setlength\tabcolsep{.4em}
% \begin{table}[t]
% \centering
% \footnotesize

\begin{tabular}{ ccccc }
\toprule
 \multirow{2}{*}{Window size}  & \multicolumn{3}{c}{DAVIS First} \\
\cmidrule(lr){2-4}
% \cmidrule(lr){5-6}
            & AJ            & \! \! \! $< \delta^x_{avg}$\! & OA    \\
\midrule

4         &  56.7  & 73.1 &  86.5  \\
\textbf{8}         &  \textbf{62.2}          & \textbf{75.7}  & \textbf{89.3} \\
16 &  61.1 & 75.5 & 88.4 \\
\bottomrule
\end{tabular}
\caption{  \textbf{Training sliding window size}. We train and evaluate \method with varying window lengths. As the training data only contains sequences of length 24, the model does not benefit from training with a bigger sliding window.}\label{tab:tab_train_sliding_window_len_ablation}
% \end{table}

    \end{minipage}%
    \hspace{0.02\linewidth}
    \begin{minipage}{0.5\linewidth}
        \centering
        \setlength\tabcolsep{1pt}
        \begin{tabular}{ cccc }
\toprule
 \multirow{2}{*}{Window size} & \multicolumn{3}{c}{DAVIS First} \\
\cmidrule(lr){2-4}
                              & AJ                              &  \! \! \! $< \delta^x_{avg}$\! & OA            \\
\midrule
4                             & 54.9                            &  70.8                          & 83.2          \\
\textbf{8}                    & \textbf{62.2}                   &  \textbf{75.7}                 & \textbf{89.3} \\
16                            & 54.6                            &  68.8                          & 82.9          \\
\bottomrule
\end{tabular}
\caption{\textbf{Inference sliding window size}. The model is trained with sliding window size $T=8$.
The same window size gives the best results.}%
\label{tab:tab_inf_sliding_window_len_ablation}
\vspace{2em}

    \end{minipage}
\end{table}

\subsection{Efficiency}

In \cref{fig:efficiency} we evaluate the efficiency of PIPs++~\cite{zheng2023pointodyssey}, TAPIR~\cite{doersch2023tapir} and \method on a 50-frame video of resolution $256\times256$ by running them on an A40 GPU\@.
We track points sampled on a regular grid at the first frame of the video and report the time it takes to process the entire video sequence.
CoTracker is slower than TAPIR but faster than PIPs++ (while achieving better accuracy than both).

\subsection{Evaluation on TAP-Vid-Kinetics}

In \cref{tab_kinetics} we evaluate \method on TAP-Vid-Kinetics, a dataset of 1144 videos of approximately 250 frames each from Kinetics~\cite{carreira2017quo}.
Some of these videos are discontinuous, i.e., they are composed of continuous video chunks.
\method is designed for continuous videos, and TAPIR contains a matching stage inspired by TAP-Net~\cite{tapvid}, which can help with such combined video chunks.
This is why the performance gap between TAPIR and \method for this benchmark is smaller than for other benchmarks.

\begin{table}[tb]
    \begin{minipage}{.47\linewidth}
        \centering
        \setlength\tabcolsep{0.9pt}
        \begin{tabular}{cccc}
\toprule
 \multirow{2}{*}{Stride} &  \multicolumn{3}{c}{Kinetics First} \\
\cmidrule(lr){2-4}
                         &  AJ                                 &  $\delta^{vis}_{avg}$ &  OA           \\
\midrule
PIPs++                   &  ---                                &  58.5                 &  ---          \\
TAPIR                    &  \textbf{49.6}                      &  64.2                 &  85.0         \\
CoTracker (Ours)         &  48.8                               &  \textbf{64.5}        &  \textbf{85.8}
                         \\ \bottomrule
\end{tabular}
\caption{
\textbf{Evaluation on TAP-Vid-Kinetics}. Kinetics contains videos with multiple shots, which does not satisfy the assumptions of most trackers, including CoTracker. The results of CoTracker and TAPIR are comparable, even though TAPIR is an offline method with a specifically designed matching module.
}\label{tab_kinetics}
\vspace{2em}

    \end{minipage}%
    \hspace{0.02\linewidth}
    \begin{minipage}{0.5\linewidth}
        \centering
        \setlength\tabcolsep{1pt}
        \begin{tabular}{cccc}
\toprule
\multirow{2}{*}{Stride} & \multicolumn{3}{c}{DAVIS First} \\
\cmidrule(lr){2-4}
                        & AJ            & $\delta^{vis}_{avg}$ &  OA            \\
\midrule
 8                      & 52.5          & 68.6                 &  85.7          \\

\textbf{4}              & \textbf{62.2} & \textbf{75.7}        &  \textbf{89.3} \\
\bottomrule
\end{tabular}
\caption{
\textbf{Feature stride ablation}. We ablate the stride of the Feature CNN\@.
Higher resolution features help to make much more accurate predictions.}%
\label{tab:tab_model_stride_ablation}
\vspace{5.5em}

    \end{minipage}
\end{table}
\begin{figure}[tb]
    \begin{minipage}{.47\linewidth}
        \centering
        \setlength\tabcolsep{0.9pt}
       \includegraphics[width=\textwidth]{figures/plot_speed.pdf}
\caption{\textbf{Efficiency}. We track N points sampled on the first frame of a 50-frame video of resolution $256 \times 256$ using a NVIDIA A40 48GB GPU, and report the time it takes the method to process the video.
}%
\label{fig:efficiency}


    \end{minipage}%
    \hspace{0.02\linewidth}
    \begin{minipage}{0.5\linewidth}
        \centering
        \setlength\tabcolsep{1pt}
        
\includegraphics[width=\textwidth]{figures/plot_ablate_updates.pdf}
\caption{\textbf{Inference iterative updates.} We ablate the number of iterative updates $M$ during inference. The network is trained with $M=4$.}%
\label{fig:plot_ablate_updates}
\vspace{2.5em}


    \end{minipage}
\end{figure}
\subsection{Implementation Details}

In this section, we complete the description of implementation details from the main paper.
We will release the code and models with the paper.

\paragraph{Feature CNN.}

Given a sequence of $T'$ frames with a resolution of 384$\times$512, we compute features for each frame using a 2-dimensional CNN\@.
This CNN downsamples the input image by a factor of 4 (feature stride) and outputs features with 128 channels.
Our CNN is the same as the one used in PIPs~\cite{pips}.
It consists of one $7\times 7$ convolution with a stride of 2, eight residual blocks with $3\times 3$ kernels and instance normalization, and two final convolutions with $3\times 3$ and $1 \times 1$ kernels.
In \cref{tab:tab_model_stride_ablation}, we train \method with two different feature downsampling factors (strides).

\paragraph{Sliding windows.}

When passing information from one sliding window to the next, we concatenate binary masks of shape $(N, T)$ with visibility logits that indicate where the model needs to make predictions.
For example, masks in the first sliding window would be equal to 1 from the frame where we start tracking the point, and 0 before that.
Masks for all the subsequent sliding windows will be equal to 1 for the first overlapping $T/2$ frames, and 0 for the remaining $T/2$.
During training, tracking starts either from the first or from a random frame where the point is visible.
If a point is not yet visible in the current sliding window, it will be masked out during the application of cross-track attention.

\paragraph{Iterative updates.}

We train the model with $M=4$ iterative updates and evaluate it with $M=6$.
This setting provides a reasonable trade-off between speed and accuracy, as evaluated in \cref{fig:plot_ablate_updates}. Interestingly, the performance remains stable for 4--8 iterative updates and begins to degrade slowly thereafter.

\paragraph{Training.}

\method is trained with a batch size of 32, distributed across 32 GPUs.
After applying data augmentations, we randomly sample 768 trajectories for each batch, with points visible either in the first or in the middle frame.
We train the model for 50,000 iterations with a learning rate of $5e^{-4}$ and a linear 1-cycle~\cite{smith2019super} learning rate schedule, using the AdamW~\cite{loshchilov2017decoupled} optimizer.

\paragraph{Augmentations.}

During training, we employ Color Jitter and Gaussian Blur to introduce color and blur variations.
We augment occlusions by either coloring a randomly chosen rectangular patch with its mean color or replacing it with another patch from the same image.
Random scaling across height and width is applied to each frame to add diversity.

\subsection{Broader societal impact}

Motion estimation, whether in the guise of point tracking or optical flow, is a fundamental, low-level computer vision task.
Many other tasks in computer vision build on it, from 3D reconstruction to video object segmentation.
Ultimately, motion estimation algorithms are important components of a very large number of applications of computer vision in many different areas.
Point tracking has no \emph{direct} societal impact; however, positive or negative effects on society can materialise through its use in other algorithms, depending on the final application.

\subsection{Technical note}

We work out the coordinate mappings between layers in a neural network.
We assume that layers are 1D, as the generalization to several dimensions is immediate.

Consider a tensor $x \in \mathbb{R}^{W}$ with $W_1$ components indexed by $i\in \Omega = \{0,\dots,W_1-1\}$.
We also interpret each element in the tensor as a unit tile, collectively covering the interval $U = [0,W]$, and we denote via $u \in U$ the corresponding continuous coordinate.
Mapping each index to the center of the corresponding tile, we have the correspondence:
$$
\Omega \rightarrow U,
~~~
i \mapsto u(i) = i + \frac{1}{2}.
$$

We interpret the tensor $x$ as providing information about image points $u \in U$.
By reading off element $x_i$, we obtain information about location $u(i)$.

\paragraph{Layers and transformations.}

Next consider two tensors $x_1 = \phi(x_2)$ related by a neural network layer $\phi$.
The layer establishes a mapping $i_1(i_2)$, or equivalently $u_1(u_2)$, between indices/coordinates of the two tensors.
We write $i_1 = \alpha_{12} i_2 + \beta_{12}$, so that the following commutative diagram holds:
\begin{center}
\begin{tikzcd}
u_2 \in U_2
\arrow{d}[swap]{\alpha_{12} u_2 + \beta_{12} + \frac{1 - \alpha_{12}}{2}}  &
i_2 \in \Omega_2
\arrow{d}{\alpha_{12} i_2 + \beta_{12}}
\arrow{l}{+\frac{1}{2}} \\
u_1 \in U_1 & i_1 \in \Omega_1
\arrow{l}[swap]{+\frac{1}{2}}
\end{tikzcd}
,
\end{center}
If we have chains of layers $x_0 \mapsto x_1 \mapsto x_2$, we can chain the corresponding diagrams:
\begin{center}
\begin{tikzcd}
u_2 \in U_2
\arrow{d}[swap]{\alpha_{12} u_2 + \beta_{12} + \frac{1 - \alpha_{12}}{2}}
&
i_2 \in \Omega_2
\arrow{d}{\alpha_{12} i_2 + \beta_{12}}
\arrow{l}{+\frac{1}{2}}
\\
u_1 \in U_1
\arrow{d}[swap]{\alpha_{1} u_1 + \beta_{1} + \frac{1 - \alpha_{1}}{2}}
& i_1 \in \Omega_1
\arrow{l}[swap]{+\frac{1}{2}}
\arrow{d}{\alpha_{1} i_1 + \beta_{1}}
\\
u_0 \in U_0
& i_0 \in \Omega_0
\arrow{l}[swap]{+\frac{1}{2}}
\end{tikzcd}
,
\end{center}
where for simplicity we denote $\alpha_{0i} = \alpha_i$ and $\beta_{0i} = \beta_i$.
With this, we find the recurrent equation:
$$
\boxed{
\alpha_n = \alpha_{n-1} \alpha_{n-1,n},
~~~
\beta_n = \beta_{n-1} + \alpha_{n-1} \beta_{n-1,n}
}
$$
so that we can always refer any layer back to the input image as follows:
$$
i_0 = \alpha_n i_n + \beta_n,~~~
u_0 = \alpha_n u_n + \beta_n + \frac{1-\alpha_{n}}{2}.
$$

Finally, we consider a sampling layer, implemented using, e.g., bilinear interpolation.
Applied to an input tensor $x_1$, we can write the latter as reading off values $x_2(u_2)$ at designated coordinates $u_2$.
This is similar to computing an output tensor $x_2$ as before, but with the ability of reading off values at arbitrary coordinates $u_2$ instead of fixed indices $i_2$.
This is written down as a function
$
x_2(u_2) =\phi(x_1, u_2).
$
Just like before, we write correspondences:
$$
i_1 = \alpha_{12} i_2 + \beta_{12},~~~
u_1 = \alpha_{12} u_2 + \beta_{12} + \frac{1-\alpha_{12}}{2}
$$
although, obviously, the index $i_2 = u_2 - 1/2$ is ``virtual''.

\paragraph{Filter-like layers.}

Next consider two tensors $x_1 = \phi(x_2)$ where $\phi$ is a neural network layer.
Assume that the layer is a filter with filter size $F$, padding $P$ and stride $S$ (almost all layers in a neural network are of this kind, often special cases of this).
Pixel $i_2$ is obtained by combinations of the values of pixels $i_1 \in \{ i_2 S - P,\dots,i_2 S - P + F-1 \}$ of the input.
Hence, we can assume that the pixel of index $i_2$ corresponds to the center of this range, which has coordinate:
$$
u_1(i_1)
= \frac{
\left(u_1(i_2 S - P) - \frac{1}{2}\right) +
\left(u_1(i_2 S - P + F-1) + \frac{1}{2}\right)
}{2}
= i_2 S - P + \frac{F-1}{2} + \frac{1}{2}.
$$
Hence:
$$
\boxed{\alpha_{12} = S,~~~
\beta_{12} = \frac{F-1}{2} - P}.
$$
Note that we can make $\beta_{12} = 0$ by choosing $P = (F-1)/2$ (which requires odd-sized filters).
By using the recurrence equation above, with a chain of such layers we obtain:
$$
\alpha_n = \prod_{i=1}^n S_n,~~~
\beta_n = 0.
$$

\paragraph{Interpolation layers.}

Next, we consider the interpolation layer $x_2 = \phi(x_1)$.
With the \texttt{align\_corners} option set to \texttt{True}, pixels $i_1=0,W_1-1$ are mapped to pixels $i_2=0,W_2-1$.
Hence:
$$
i_1
= i_2 \frac{W_1-1}{W_2-1}
= \alpha_{12}i_2 + \beta_{12}.
$$
Hence we have:
$$
\alpha_{12} = \frac{W_1-1}{W_2-1},~~~
\beta_{12} = 0.
$$
With the \texttt{align\_corners} option set to \texttt{False}, coordinates $u_1=0,W_1$ are mapped to coordinates $u_2=0,W_2$.
Hence:
$$
u_2
= u_1 \frac{W_1}{W_2}
= \alpha_{12} u_1
+ \beta_{12}
+ \frac{1 - \alpha_{12}}{2}.
$$
Hence:
$$
\alpha_{12} = \frac{W_1}{W_2},~~~
\beta_{12} = - \frac{1 - \alpha_{12}}{2}
= \frac{W_1-W_2}{2W_2}.
$$

\paragraph{Sampling layers.}

Next, we figure out the sampling layer $x_2(u_2) = \phi(x_1,u_2)$.
When we use bilinear sampling to sample features from $x_1$, we use the \texttt{grid\_sample} operator with \texttt{align\_corners} set to either \texttt{True} or \texttt{False}.

We can model this as introducing a coordinate $u_2 \in U_2 = [-1,1]$ with the following mappings.
If \texttt{align\_corners} is \texttt{False}, we align the extrema $-1,1$ of the output range to the edges of the input image, so that:
$$
u_1(u_2)
=
\frac{W_1}{2} (u_2 + 1)
=
\alpha_{12} u_1
+ \beta_{12}
+ \frac{1 - \alpha_{12}}{2}.
$$
i.e.,
$$
\alpha_{12} = \frac{W_1}{2},~~~
\beta_{12} = \frac{W_1}{2} - \frac{1 - \alpha_{12}}{2} = - \frac{1}{2} - \frac{W_1}{4}.
$$
If \texttt{align\_corners} is \texttt{True}, we align the extrema of the output range $-1, 1$ with the centers of the input edge pixels, so that:
$$
i_1(u_2)
=
\frac{W_1-1}{2} (u_2 + 1)
=
\alpha_{12}u_2 + \beta_{12},
$$
i.e.,
$$
\alpha_{12} = \frac{W_1 - 1}{2},~~~
\beta_{12} = \frac{W_1 - 1}{2}.
$$

\paragraph{Sampling layer, alt.~conventions.}

In practice, we redefine the mapping layers in such a way that, if \texttt{align\_corners} is \texttt{False}, then $u_2 \in U_2 = [0, W_1]$ and the extrema $0, W_1$ align edges of the input image, so that:
$$
u_1(u_2)
=
u_2
=
\alpha_{12} u_1
+ \beta_{12}
+ \frac{1 - \alpha_{12}}{2}.
$$
i.e.,
$$
\alpha_{12} = 1,~~~
\beta_{12} = 0,
$$
which means:
$$
u_1(u_2) = u_2, ~~~
i_1(i_2) = i_2.
$$
With \texttt{align\_corners} is \texttt{True}, then $u_2 \in U_2 = [0, W_1-1]$ and we align the extrema $0, W_1-1$ to the center of the edge pixels, so that:
$$
i_1(u_2)
=
u_2
=
\alpha_{12}u_2 + \beta_{12} - \frac{1-\alpha_{12}}{2} - \frac{1}{2},
$$
i.e.,
$$
\alpha_{12} = 1,~~~
\beta_{12} = \frac{1}{2},
$$
which means:
$$
u_1(u_2) = u_2 + \frac{1}{2},~~~
i_1(i_2) = i_1 + \frac{1}{2}.
$$